\section{Problem and Threat Model}
\label{sec:problem}

\subsection{System and Trust Model}

We consider a blockchain-based healthcare data-sharing system in which
access to protected resources is governed by provider-defined consent
policies. The system involves five roles: a data provider or data owner, a data
requester, a Trusted Proof Service (TPS), an on-chain access-control manager
(ACM), and a resource access layer. The provider defines a consent policy $P$
for a protected resource and registers an on-chain commitment $h$ to the encoded
policy. The requester seeks access by submitting policy-relevant
authorization evidence, such as attributes, credentials, and access context.
The TPS constructs the private witness and generates a zero-knowledge proof
of policy compliance. The ACM verifies the proof on-chain and records the authorization
outcome, while the resource access layer grants access according to the
recorded decision.

TPS-ZKPI relies on the following assumptions. First, policies are provider-controlled:
the provider is responsible for defining the applicable consent policy and committing
the corresponding policy instance on-chain. Second, the TPS is trusted to preserve
the confidentiality of policy-bound proving material and requester-submitted
authorization evidence, and to follow the prescribed proof-generation
workflow. Third, the commitment scheme is assumed to be hiding and binding, and the zero-knowledge proof system is assumed to satisfy completeness, soundness, and zero knowledge with respect to the implemented policy-compliance relation. Credential
issuance, revocation, and real-world identity verification are treated as
external functions.

The evaluated instantiation uses a provider-operated TPS. This choice avoids exposing
the policy encoding or commitment-opening randomness to requesters and keeps public
verification stable, but it also concentrates trust at the provider-side proof service:
a compromised or colluding TPS can see both the policy and requester evidence, and a
provider-operated TPS can reintroduce fabrication and stonewalling risks similar to a
provider-prove design. The commitment and zero-knowledge proof machinery mainly buys
public verifiability and correctness against a TPS that submits an invalid proof; it
does not provide confidentiality against a compromised TPS. Reduced-trust alternatives
include an independent proof service, which becomes a high-value honeypot, a TEE-backed
proof service, which makes the trust assumption hardware-attested, or MPC/collaborative
SNARK proof generation, where no single party sees both policy material and requester
evidence.

\subsection{Public-Chain Adversary}

The primary adversary is a passive public-chain observer. The adversary can
inspect smart-contract code and monitor public blockchain artifacts across
multiple authorization sessions, including policy commitments, submitted
proofs, public inputs, emitted events, calldata, timestamps, and gas usage.
The adversary does not observe provider-to-TPS policy material, commitment-opening
randomness, requester-submitted authorization evidence, credential witnesses,
or TPS-side witness construction.

Each policy-enforcement instance corresponds to an authorization session, in
which a requester seeks access to a protected resource under a committed policy.
For each such session, the adversary observes a public view
\begin{equation}
  view_{\mathsf{pub}}= (h, \pi, id_{req}, id_{res}, nonce, pp, meta),
\end{equation}
where $h$ is the on-chain commitment to the hidden policy, $\pi$ is the zero-knowledge
proof submitted for on-chain verification, $id_{req}$, $id_{res}$, and $nonce$ are the
public requester, resource, and session-binding values, $pp$ denotes other fixed-format
public parameters or public handles, and $meta$ denotes transaction-level metadata
exposed by the enforcement layer, such as calldata footprint, event structure,
public-input layout, timing, and gas usage. The adversary's goal is to infer hidden
information about the committed policy or the authorization interaction from
this public view.

In addition to the public-chain observer, we treat the requester as untrusted
with respect to policy confidentiality. The requester may initiate access
requests, observe its own submitted evidence and authorization outcome, and
inspect the public transcript. However, the requester does not receive the
encoded policy, commitment-opening randomness, or policy-bound proving
material, which remain within the provider-to-TPS proof-generation boundary.

Compromised TPS behavior, adaptive requester probing, and off-chain side channels
are non-goals for the current model. The SNARK checks internal consistency of the
witness it is given; it does not prove that the private attribute vector is the
requester's real-world attribute set. Binding a proof to $id_{req}$ ties the proof to
an identifier and session context, but without an issuer-verification layer it does
not by itself bind every private attribute value to an external source of truth.

\subsection{Privacy Objectives}

We define policy privacy through three attack-oriented objectives.

\textit{Policy Content Hiding (PCH).} PCH requires that public-chain artifacts
should not reveal the semantic content of the hidden consent policy, such as
data-use terms, disease categories, research purposes, institutional restrictions,
or ethics-related requirements.

\textit{Policy Indistinguishability (PI).} PI requires that different candidate
policies should not become distinguishable from their public verification artifacts
when they are compiled into the same fixed public interface. This captures
whether two policies can be told apart from commitments, proofs, public
parameters, or metadata even when their direct representations are hidden.

\textit{Structure and Parameter Privacy (SPP).} SPP requires that public-chain
artifacts should not reveal structural or parametric properties of the
hidden policy, such as logical structure, nesting depth, predicate count,
set size, or threshold-like parameters.

\subsection{Problem Statement}

Given a provider-defined consent policy $P$, requester authorization evidence,
an untrusted public blockchain, and the stated TPS trust assumption, the problem is
to enforce access to a protected healthcare resource while satisfying three requirements. First, authorization
must be publicly verifiable: the ACM should be able to verify that the
requester satisfies the committed policy without learning or executing the
policy in plaintext. Second, policy-dependent information should remain
hidden from public-chain observers, including policy content, policy
distinguishability, structural properties, requester evidence, and metadata-level
enforcement fingerprints. Third, authorization outcomes should remain auditable,
so that successful access decisions can be recorded and checked without
exposing the hidden policy logic.
